\begin{ex}
Phân số $\frac{5}{9}$ có tử số là bao nhiêu?
\shortans{5}
\loigiai{
Tử số là 5.
}
\end{ex}

\begin{ex}
Phân số $\frac{5}{9}$ có mẫu số là bao nhiêu?
\shortans{9}
\loigiai{
Mẫu số là 9.
}
\end{ex}

\begin{ex}
Rút gọn phân số $\frac{12}{16}$ về tối giản.
\shortans{3/4}
\loigiai{
$\frac{3}{4}$.
}
\end{ex}

\begin{ex}
Rút gọn phân số $\frac{20}{35}$ về tối giản.
\shortans{4/7}
\loigiai{
$\frac{4}{7}$.
}
\end{ex}

\begin{ex}
Rút gọn phân số $\frac{45}{60}$ về tối giản.
\shortans{3/4}
\loigiai{
$\frac{3}{4}$.
}
\end{ex}

\begin{ex}
Quy đồng mẫu số $\frac{2}{3}$ và $\frac{3}{5}$ với mẫu số chung 15, phân số $\frac{2}{3}$ trở thành phân số nào?
\shortans{10/15}
\loigiai{
$\frac{10}{15}$.
}
\end{ex}

\begin{ex}
Phân số $\frac{3}{5}$ trở thành phân số nào trong bài toán trên?
\shortans{9/15}
\loigiai{
$\frac{9}{15}$.
}
\end{ex}

\begin{ex}
So sánh $\frac{5}{8}$ và $\frac{7}{8}$: phân số nào lớn hơn?
\shortans{7/8}
\loigiai{
$\frac{7}{8}$.
}
\end{ex}

\begin{ex}
So sánh $\frac{3}{4}$ và $\frac{3}{7}$: phân số nào lớn hơn?
\shortans{3/4}
\loigiai{
$\frac{3}{4}$.
}
\end{ex}

\begin{ex}
So sánh $\frac{4}{5}$ và 1: phân số $\frac{4}{5}$ bé hơn hay lớn hơn 1? (Bé hơn ghi 0, lớn hơn ghi 1).
\shortans{0}
\loigiai{
Bé hơn 1 (0).
}
\end{ex}

\begin{ex}
Phân số $\frac{9}{8}$ lớn hơn hay bé hơn 1? (Lớn hơn ghi 1, bé hơn ghi 0).
\shortans{1}
\loigiai{
Lớn hơn 1 (1).
}
\end{ex}

\begin{ex}
Viết số 6 dưới dạng phân số có mẫu số bằng 1.
\shortans{6/1}
\loigiai{
$\frac{6}{1}$.
}
\end{ex}

\begin{ex}
Viết số 1 dưới dạng phân số có tử số bằng 15.
\shortans{15/15}
\loigiai{
$\frac{15}{15}$.
}
\end{ex}

\begin{ex}
Tìm $x$: $\frac{x}{24} = \frac{3}{4}$.
\shortans{18}
\loigiai{
$x = 3 \times 6 = 18$.
}
\end{ex}

\begin{ex}
Tìm $y$: $\frac{5}{y} = \frac{15}{21}$.
\shortans{7}
\loigiai{
$y = 21 : 3 = 7$.
}
\end{ex}

\begin{ex}
Trong các phân số $\frac{2}{3}, \frac{4}{6}, \frac{7}{8}, \frac{9}{12}$, có bao nhiêu phân số tối giản?
\shortans{2}
\loigiai{
$\frac{2}{3}$ và $\frac{7}{8}$ (có 2 phân số).
}
\end{ex}

\begin{ex}
Phân số chỉ phần đã tô màu của hình chia 8 phần bằng nhau, tô 5 phần là bao nhiêu?
\shortans{5/8}
\loigiai{
$\frac{5}{8}$.
}
\end{ex}

\begin{ex}
Phần chưa tô màu trong bài toán trên là bao nhiêu?
\shortans{3/8}
\loigiai{
$\frac{3}{8}$.
}
\end{ex}

\begin{ex}
Thương của phép chia $7 : 11$ viết dưới dạng phân số là bao nhiêu?
\shortans{7/11}
\loigiai{
$\frac{7}{11}$.
}
\end{ex}

\begin{ex}
Mẫu số chung nhỏ nhất của $\frac{1}{4}$ và $\frac{1}{6}$ là bao nhiêu?
\shortans{12}
\loigiai{
12.
}
\end{ex}

\begin{ex}
Rút gọn phân số $\frac{30}{40}$ về tối giản.
\shortans{3/4}
\loigiai{
$\frac{3}{4}$.
}
\end{ex}

\begin{ex}
Phân số $\frac{14}{28}$ rút gọn về tối giản bằng bao nhiêu?
\shortans{1/2}
\loigiai{
$\frac{1}{2}$.
}
\end{ex}

\begin{ex}
So sánh $\frac{2}{3}$ và $\frac{5}{6}$: phân số nào lớn hơn?
\shortans{5/6}
\loigiai{
$\frac{4}{6} < \frac{5}{6}$.
}
\end{ex}

\begin{ex}
Trong các phân số $\frac{1}{2}, \frac{2}{3}, \frac{3}{4}$, phân số lớn nhất là phân số nào?
\shortans{3/4}
\loigiai{
$\frac{3}{4}$.
}
\end{ex}

\begin{ex}
Phân số bé nhất trong 3 phân số trên là phân số nào?
\shortans{1/2}
\loigiai{
$\frac{1}{2}$.
}
\end{ex}
